% GENERATED FILE. Edit canonical lessons, then run npm run generate:books.
% canonical-source-hash: fnv1a64:e028d07b5ae38e02
% canonical-lessons: LA-C230-expergiscor, LA-C230-comitor, LA-C230-addo, LA-C230-nuntio, LA-C230-appareo

\chapter{To wake up, To accompany, To add, To announce, To appear}
\label{ch:la-to-wake-up-to-accompany-to-add-to-announce-to-appear}
\hlchaptermodality{230}

\begin{chapteropening}
\textbf{By the end of this chapter:} \emph{I can say to wake up, to accompany, to add, to announce and to appear.}

Complete the last lesson of chapter 230: I can say to wake up, to accompany, to add, to announce and to appear.
\end{chapteropening}

\section[expergīscor, expergīscī]{expergīscor, expergīscī — to wake up}

\label{lesson:LA-C230-expergiscor}



\begin{quote}
Before the new one: say the Latin for to hold off, then the Latin for to remove.
\end{quote}

\subsection*{You'll want to know: expergīscor, expergīscī}
\textbf{expergīscor, expergīscī} — \textquotedblleft{}to wake up\textquotedblright{}.

\begin{cousinweb}
\textbf{expergīscor} is built on \emph{expergō}, \textquotedblleft{}to wake up\textquotedblright{} (Lewis \& Short), with \emph{-scor}, the ending of verbs of becoming; \emph{expergō} is often explained as \emph{ex}, \textquotedblleft{}out\textquotedblright{}, with \emph{pergō}, \textquotedblleft{}to go on\textquotedblright{}.
\end{cousinweb}

\begin{grammarlens}[title={What you've built}]
One more word for this chapter.
\end{grammarlens}

\subsection*{Your turn}
\begin{itemize}
\raggedright
  \item \emph{Say it:} \emph{expergīscor, expergīscī}
  \item \emph{Say it:} \emph{expergīscor, expergīscī}, once more
  \item \emph{Say it:} say \emph{expergīscor, expergīscī}
  \item \emph{Recall:} say the Latin for to hold off, then the Latin for to remove, then say \emph{expergīscor, expergīscī} again
\end{itemize}

\begin{tcolorbox}[breakable,colback=teal!4,colframe=teal!35!black,title={Before you move on}]
What does expergīscor, expergīscī mean? (\textquotedblleft{}To wake up\textquotedblright{}.) Say it once more.
\end{tcolorbox}
\section[comitor, comitārī]{comitor, comitārī — to accompany}

\label{lesson:LA-C230-comitor}



\begin{quote}
Before the new one: say the Latin for to remove, then the Latin for to wake up.
\end{quote}

\subsection*{You'll want to know: comitor, comitārī}
\textbf{comitor, comitārī} — \textquotedblleft{}to accompany\textquotedblright{}.

\begin{cousinweb}
\textbf{comitor} is built on \emph{comes}, \emph{comitis}, \textquotedblleft{}a companion\textquotedblright{} (Lewis \& Short): to be someone's companion. \emph{comes} is \emph{con-}, \textquotedblleft{}together\textquotedblright{}, with the root of \emph{eō}, \textquotedblleft{}to go\textquotedblright{}: one who goes along.
\end{cousinweb}

\begin{grammarlens}[title={What you've built}]
One more word for this chapter.
\end{grammarlens}

\subsection*{Your turn}
\begin{itemize}
\raggedright
  \item \emph{Say it:} \emph{comitor, comitārī}
  \item \emph{Say it:} \emph{comitor, comitārī}, once more
  \item \emph{Say it:} say \emph{comitor, comitārī}
  \item \emph{Recall:} say the Latin for to remove, then the Latin for to wake up, then say \emph{comitor, comitārī} again
\end{itemize}

\begin{tcolorbox}[breakable,colback=teal!4,colframe=teal!35!black,title={Before you move on}]
What does comitor, comitārī mean? (\textquotedblleft{}To accompany\textquotedblright{}.) Say it once more.
\end{tcolorbox}
\section[addō, addere]{addō, addere — to add}

\label{lesson:LA-C230-addo}



\begin{quote}
Before the new one: say the Latin for to wake up, then the Latin for to accompany.
\end{quote}

\subsection*{You'll want to know: addō, addere}
\textbf{addō, addere} — \textquotedblleft{}to add\textquotedblright{}.

\begin{cousinweb}
\textbf{addō} is \emph{ad}, \textquotedblleft{}to\textquotedblright{}, with \emph{dō}, \textquotedblleft{}to put, to give\textquotedblright{} (Lewis \& Short): to put to.
\end{cousinweb}

\begin{grammarlens}[title={What you've built}]
One more word for this chapter.
\end{grammarlens}

\subsection*{Your turn}
\begin{itemize}
\raggedright
  \item \emph{Say it:} \emph{addō, addere}
  \item \emph{Say it:} \emph{addō, addere}, once more
  \item \emph{Say it:} say \emph{addō, addere}
  \item \emph{Recall:} say the Latin for to wake up, then the Latin for to accompany, then say \emph{addō, addere} again
\end{itemize}

\begin{tcolorbox}[breakable,colback=teal!4,colframe=teal!35!black,title={Before you move on}]
What does addō, addere mean? (\textquotedblleft{}To add\textquotedblright{}.) Say it once more.
\end{tcolorbox}
\section[nūntiō, nūntiāre]{nūntiō, nūntiāre — to announce}

\label{lesson:LA-C230-nuntio}



\begin{quote}
Before the new one: say the Latin for to accompany, then the Latin for to add.
\end{quote}

\subsection*{You'll want to know: nūntiō, nūntiāre}
\textbf{nūntiō, nūntiāre} — \textquotedblleft{}to announce\textquotedblright{}.

\begin{cousinweb}
\textbf{nūntiō} is the verb made from \emph{nūntius}, \textquotedblleft{}a messenger, news\textquotedblright{} (Lewis \& Short), whose origin is uncertain.
\end{cousinweb}

\begin{grammarlens}[title={What you've built}]
One more word for this chapter.
\end{grammarlens}

\subsection*{Your turn}
\begin{itemize}
\raggedright
  \item \emph{Say it:} \emph{nūntiō, nūntiāre}
  \item \emph{Say it:} \emph{nūntiō, nūntiāre}, once more
  \item \emph{Say it:} say \emph{nūntiō, nūntiāre}
  \item \emph{Recall:} say the Latin for to accompany, then the Latin for to add, then say \emph{nūntiō, nūntiāre} again
\end{itemize}

\begin{tcolorbox}[breakable,colback=teal!4,colframe=teal!35!black,title={Before you move on}]
What does nūntiō, nūntiāre mean? (\textquotedblleft{}To announce\textquotedblright{}.) Say it once more.
\end{tcolorbox}
\section[appāreō, appārēre]{appāreō, appārēre — to appear}

\label{lesson:LA-C230-appareo}



\begin{quote}
Before the new one: say the Latin for to add, then the Latin for to announce.
\end{quote}

\subsection*{You'll want to know: appāreō, appārēre}
\textbf{appāreō, appārēre} — \textquotedblleft{}to appear\textquotedblright{}.

\begin{cousinweb}
\textbf{appāreō} is \emph{ad}, \textquotedblleft{}to\textquotedblright{}, with \emph{pāreō}, \textquotedblleft{}to come into view\textquotedblright{} (EtymDB 2.1). English \emph{appear} comes from it through French.
\end{cousinweb}

\begin{grammarlens}[title={What you've built}]
One more word for this chapter.
\end{grammarlens}

\subsection*{Your turn}
\begin{itemize}
\raggedright
  \item \emph{Say it:} \emph{appāreō, appārēre}
  \item \emph{Say it:} \emph{appāreō, appārēre}, once more
  \item \emph{Say it:} say \emph{appāreō, appārēre}
  \item \emph{Recall:} say the Latin for to add, then the Latin for to announce, then say \emph{appāreō, appārēre} again
\end{itemize}

\begin{tcolorbox}[breakable,colback=teal!4,colframe=teal!35!black,title={Before you move on}]
What does appāreō, appārēre mean? (\textquotedblleft{}To appear\textquotedblright{}.) Say it once more.
\end{tcolorbox}
